\chapter{도파민이 실제로 말하는 것}
% full version: chapters/07_dofaalt.tex

지난 장에서 도파민은 단순했다. 숫자 하나, “예상보다 좋은가 나쁜가”. 이것은 1차 근사이고 잘 들어맞는다. 그러나 최근 몇 년의 실험은 이 “전선”으로 숫자 하나보다 많은 것이 흐른다는 것을 보여 주었다. 이 장은 그림이 더 복잡해지는 곳, 그리고 학자들이 아직 논쟁하는 곳을 다룬다.

\section*{낙관주의자와 비관주의자}

도파민 뉴런은 모두 같지 않다. 어떤 뉴런은 기분 나쁜 뜻밖의 일보다 기분 좋은 뜻밖의 일에 더 세게 반응한다. 낙관주의자이다. 반대인 뉴런은 비관주의자이다. 뉴런마다 제 규칙대로 배우면, 낙관주의자는 가장 좋은 경우의 보상을, 비관주의자는 가장 나쁜 경우의 보상을 배운다. 이들은 함께 평균 기대 하나가 아니라 가능한 결과의 분포 전체를 기록한다. 우승 후보만이 아니라 모든 참가자의 승률을 아는 도박 중개인과 같다. 이 생각은 실험과 들어맞는다. 뇌가 왜 분포 전체를 필요로 하는지는 아직 가설이다.

\pencilfig{7}{\pencilart{7}}{도파민 뉴런은 저마다 가능한 보상 곡선 위의 한 점을 맡는다. 함께 모이면 곡선 전체를 기억한다.}

\section*{“더 좋다”만이 아니라 “중요하다”도}

도파민 뉴런 일부는 불쾌하지만 중요한 사건, 이를테면 큰 소리나 위험에도 분출한다. 도파민 네트워크로는 뜻밖의 일의 부호만이 아니라 그 중요도, 즉 “주의하라”는 신호도 흐르는 것 같다.

\section*{전선 하나에 두 신호}

빠른 분출은 가르친다. 그런데 몇 분에 걸쳐 변하는 도파민의 느린 평균 수준은 다른 것을 말하는 듯하다. 지금 주변이 얼마나 풍요로운가이다. 주변에 보상이 많으면 시간이 귀하고, 동물은 더 빨리 움직인다. 실험은 도파민 수준이 행동의 활기와 관련 있음을 확인한다. 엔지니어라면 이것을 주파수 분할이라고 부를 것이다. 높은 주파수는 한 가지를, 낮은 주파수는 다른 것을 나른다.

\section*{목표 앞에서 차오르기}

동물이 보상에 다가갈 때 도파민은 분출하지 않고 서서히 오르는 경우가 많다. 한 가지 설명은 이것이 기대 그 자체, 즉 “목표가 가깝다, 힘내라”는 신호라는 것이다. 다른 설명은 뜻밖의 일이라는 생각을 유지한다. 멀리서는 동물이 상황을 흐릿하게 평가하고, 다가갈수록 점점 정확하게 평가하며, 정확해질 때마다 작은 기분 좋은 뜻밖의 일이 생긴다는 것이다. 가상 복도에서 동물을 순간적으로 목표 가까이 옮긴 멋진 실험은 두 번째 설명을 지지한다. 도파민은 뜻밖의 일에 걸맞게 분출로 응답했다.

\section*{뒤돌아보기}

과감한 대안도 있다. 도파민은 “무엇이 일어날까?”가 아니라 “무엇이 보상의 원인이었나?”에 답한다는 것이다. 미래가 아니라 과거를 본다. 두 이론이 서로 다른 예측을 내놓는 실험들은 아직 서로 어긋난다. 살아 있는 과학 논쟁이다.

\section*{숫자 대신 벡터}

끝으로, 도파민은 보상의 가치는 그대로인데 \emph{맛}이 바뀔 때도 반응하고, 보상이 전혀 없는 학습도 돕는다. 이것은 오차 하나가 아니라 오차의 묶음인 것 같다. 일어난 일의 특징 하나마다 오차 하나씩이다.

이 장의 결론: 새 이론 대부분은 단순한 그림을 뒤엎지 않고 넓힌다. 도파민은 전선 한 가닥이 아니라 몇 가닥으로 된 작은 다발이다. “뒤돌아보기”만이 단순한 그림과 진지하게 맞선다.

\fullversion{7}
